%% Lecture 17 -- Magnetization currents, the H field, and the magnetic scalar potential
\section{Lecture 17 -- Magnetization Currents, the $\Hvec$ Field, and the Magnetic Scalar Potential}\label{sec:lec17}

\begin{center}
\begin{tabular}{ll}
  \textbf{Date:}\quad 02/06/2026 & \\
\end{tabular}
\end{center}
\hrule\vspace{1.5em}

\subsection*{Overview}

Lecture 16 established that a piece of matter behaves, on macroscopic length scales,
as a continuous distribution of magnetic dipoles, summarized by the magnetization
$\bvec{M}(\br)=d\bmu_{\rm total}/dV$.  This lecture closes the magnetostatics of matter
in three steps:
\begin{itemize}
  \item identify the microscopic origin of $\bvec{M}$ (orbital currents and spin) and
        write the \emph{total} current as free plus magnetization currents;
  \item start from the dipole vector potential, superpose it over the body, and prove
        that magnetization is equivalent to a bound \emph{volume} current
        $\Jvec_M=c\,\nabla\times\bvec{M}$ plus a bound \emph{surface} current
        $\bvec{K}_M=c\,\bvec{M}\times\uvec{n}$;
  \item insert these into Amp\`ere's law, which forces the definition of the magnetic
        intensity $\Hvec=\bB-4\pi\bvec{M}$ sourced only by free currents, and -- when
        free currents are absent -- a \emph{magnetic scalar potential} obeying a
        Poisson equation identical in form to electrostatics.
\end{itemize}
The recurring theme is that everything magnetic matter does is captured by $\bvec{M}$:
once $\bvec{M}$ (or equivalently $\Jvec_M$) is given, the field is fully determined.

%------------------------------------------------------------
\subsection{Total Current in Magnetic Matter}\label{subsec:lec17-total-current}

\subsubsection*{Physical Motivation}
A magnetic field is produced by moving charge.  In a laboratory there are two very
different kinds of moving charge.  The first is the current we deliberately drive
through wires -- the current in the solenoid we use to magnetize a sample.  The second
appears only once matter is present: the perpetual microscopic currents bound inside
atoms and molecules.  In equilibrium these microscopic currents are randomly oriented
and average to zero, but an external field can order them and produce a net effect.
The macroscopic field therefore responds to the \emph{sum} of these contributions.

\dfn[dfn:lec17-current-split]{Free and Magnetization Currents}{
  The total current density measured at macroscopic (classical) resolution splits as
  \begin{equation}\label{eq:lec17-J-total}
    \boxed{\Jvec_{\rm total}=\Jvec_{\rm f}+\Jvec_M.}
  \end{equation}
  $\Jvec_{\rm f}$ are \emph{free currents}: conduction electrons travelling long
  distances through wires.  $\Jvec_M$ are \emph{magnetization} (or bound) currents:
  the coarse-grained average of the microscopic currents bound to atoms.
}
\textit{Significance:} Only $\Jvec_{\rm f}$ is under direct experimental control; the
whole physics of magnetic matter lives in how $\Jvec_M$ is generated and how it relates
to $\bvec{M}$.

\subsubsection*{Microscopic Origin of the Magnetization Current}
The measured macroscopic current is an average of a microscopic current that itself has
two pieces,
\begin{equation}\label{eq:lec17-Jm-split}
  \Jvec_m=\Jvec_o+\Jvec_s,
\end{equation}
where $\Jvec_o$ is the genuine \emph{orbital} current produced by the motion of bound
charges, and $\Jvec_s$ is an \emph{effective} current standing in for electron spin.

For the orbital part, a small volume $\Delta V$ contains a huge number of microscopic
current loops; coarse-graining replaces them by their dipole moment per unit volume.
Using the dipole-moment definition $\bmu=\tfrac{1}{2c}\int \br'\times\Jvec\,d^3x'$ from
Lecture 15, the orbital magnetization is
\begin{equation}\label{eq:lec17-Mo-def}
  \boxed{\bvec{M}_o(\br)=\frac{1}{2c\,\Delta V}\int_{\Delta V}\br'\times\Jvec_o(\br')\,d^3x'
  =\frac{d\bmu_o}{dV}.}
\end{equation}
\textit{Significance:} The complicated microscopic loop structure is compressed into a
single field, the dipole-moment density; this is the only orbital information that
survives at macroscopic distances.

\subsubsection*{The Spin Contribution}
Spin is an intrinsic angular momentum of the electron with no classical analogue: it is
\emph{not} the circulation of any charge.  Nevertheless, because orbital angular
momentum $\bvec{L}$ produces a magnetic moment $\bmu=(e/2m_ec)\bvec{L}$, and spin
$\bvec{S}$ is also an angular momentum, consistency demands that spin produce a magnetic
moment too.  We therefore model its effect by an effective current whose magnetic moment
matches that of spin.  Its magnetization is a sum over the spins per unit volume,
\begin{equation}\label{eq:lec17-Ms-def}
  \bvec{M}_s(\br)=\sum_{k}\bvec{m}_k\,\delta^{(3)}(\br-\br_k),
  \qquad
  \bvec{m}_k\propto \bvec{S}_k\sim\frac{e}{2m_ec}\,\bvec{S}_k,
\end{equation}
the sum running over all spins inside a unit volume.  Although the underlying cause of
$\bvec{M}_s$ is purely quantum, the resulting magnetic effect is measured by ordinary
classical instruments.

\nt{The magnitude of spin (the quantum number $\tfrac12$) is Lorentz invariant like a
rest mass, but spin is a \emph{vector}: its direction can change, which is exactly what
lets an external field reorient it. The total angular momentum -- and hence the total
current -- does transform between frames, while the spin magnitude does not.}

\subsubsection*{Total Magnetization}
For our (classical, phenomenological) purposes we never separate the two sources: a
classical experiment sees neither the orbital motion around an atom nor the spin
directly.  We work only with the total magnetization
\begin{equation}\label{eq:lec17-M-total}
  \boxed{\bvec{M}=\bvec{M}_o+\bvec{M}_s.}
\end{equation}
\textit{Significance:} The macroscopic theory needs only $\bvec{M}$; the distinction
between ``real'' orbital currents and ``effective'' spin currents is a microscopic
detail invisible to classical measurement.

%------------------------------------------------------------
\subsection{Vector Potential of a Magnetized Body}\label{subsec:lec17-vector-potential}

\subsubsection*{Physical Motivation}
Suppose the magnetization $\bvec{M}(\br)$ is given throughout a body.  We want the
field it produces.  Because the theory is linear, the total field is the superposition
of the fields of all the microscopic dipoles, plus (separately) the field of any free
currents, which we already know how to handle.  Each small volume $d^3r'$ carries a
dipole moment $d\bmu=\bvec{M}(\br')\,d^3r'$, and a dipole at $\br'$ contributes the
vector potential $\bmu\times(\br-\br')/|\br-\br'|^3$ from Lecture 15.

\subsubsection*{Mathematical Setup}
Summing (integrating) these contributions over the body gives the magnetization vector
potential
\begin{equation}\label{eq:lec17-Am}
  \boxed{\Avec_M(\br)
  =\int d^3r'\,
  \frac{\bvec{M}(\br')\times(\br-\br')}{|\br-\br'|^3}.}
\end{equation}
\textit{Significance:} This is the exact analogue of replacing a charge distribution by
its dipole density; it is the leading (dipole) approximation to the field of matter,
valid because classical probes always sit at distances enormous compared with atomic
scales.

\begin{figure}[h]
\centering
\begin{tikzpicture}[>=Stealth, thick, scale=1.0]
  \draw[->, gray] (-0.5,0) -- (5.2,0) node[right] {$x$};
  \draw[->, gray] (0,-0.5) -- (0,3.4) node[above] {$z$};
  \fill[black] (0,0) circle (1.5pt);
  \node[below left] at (0,0) {$O$};
  % source element
  \draw[rounded corners=2pt, fill=blue!8] (1.1,1.0) rectangle (1.7,1.55);
  \draw[->, red!70!black, line width=1.0pt] (1.4,1.27) -- (1.95,1.85)
    node[right] {$\bvec{M}(\br')\,d^3r'$};
  % observation point
  \fill[black] (4.2,2.7) circle (1.5pt);
  \node[above right] at (4.2,2.7) {obs.\ pt};
  % vectors
  \draw[->, blue!60!black, line width=0.9pt] (0,0) -- (1.4,1.27)
    node[midway, below right] {$\br'$};
  \draw[->, green!55!black, line width=0.9pt] (0,0) -- (4.2,2.7)
    node[midway, above left] {$\br$};
  \draw[->, orange!85!black, line width=0.9pt] (1.4,1.27) -- (4.2,2.7)
    node[midway, below right] {$\br-\br'$};
\end{tikzpicture}
\caption{Each volume element at $\br'$ carries a dipole $d\bmu=\bvec{M}(\br')\,d^3r'$
and contributes to the vector potential at the observation point $\br$ through the
separation vector $\br-\br'$.}
\label{fig:lec17-superposition}
\end{figure}

%------------------------------------------------------------
\subsection{Magnetization (Bound) Currents}\label{subsec:lec17-bound-currents}

\subsubsection*{Mathematical Setup}
We rewrite \eqref{eq:lec17-Am} so that it looks like the magnetostatic Poisson solution
$\Avec=\tfrac1c\int \Jvec/|\br-\br'|\,d^3r'$, and read off the effective current.  The
key observation is that the separation factor is a primed gradient:
\begin{equation}\label{eq:lec17-grad-prime}
  \nabla'\frac{1}{|\br-\br'|}=\frac{\br-\br'}{|\br-\br'|^3},
\end{equation}
where $\nabla'$ differentiates with respect to $\br'$.  Apply the product identity
$\nabla'\times(f\bvec{V})=(\nabla' f)\times\bvec{V}+f\,(\nabla'\times\bvec{V})$ with
$f=1/|\br-\br'|$ and $\bvec{V}=\bvec{M}(\br')$:
\begin{equation}\label{eq:lec17-product-identity}
  \nabla'\times\!\left(\frac{\bvec{M}}{|\br-\br'|}\right)
  =\left(\nabla'\frac{1}{|\br-\br'|}\right)\times\bvec{M}
  +\frac{\nabla'\times\bvec{M}}{|\br-\br'|}.
\end{equation}
Substituting \eqref{eq:lec17-grad-prime} and isolating the integrand of
\eqref{eq:lec17-Am},
\begin{align}
  \frac{\bvec{M}\times(\br-\br')}{|\br-\br'|^3}
  &=-\left(\frac{\br-\br'}{|\br-\br'|^3}\right)\times\bvec{M}
  \notag\\
  &=\frac{\nabla'\times\bvec{M}}{|\br-\br'|}
  -\nabla'\times\!\left(\frac{\bvec{M}}{|\br-\br'|}\right).
  \label{eq:lec17-integrand-rewrite}
\end{align}

\subsubsection*{Converting the Curl Term to a Surface Integral}
For the last term we use the curl form of Gauss's theorem,
\begin{equation}\label{eq:lec17-curl-gauss}
  \int_V (\nabla\times\bvec{A})\,d^3r=\oint_S \uvec{n}\times\bvec{A}\,d\sigma
  =\oint_S d\boldsymbol{\sigma}\times\bvec{A},
\end{equation}
with $\uvec{n}$ the outward normal of the body's surface.  Hence
\begin{equation}\label{eq:lec17-curl-to-surface}
  \int_V \nabla'\times\!\left(\frac{\bvec{M}}{|\br-\br'|}\right)d^3r'
  =\oint_S \frac{\uvec{n}\times\bvec{M}}{|\br-\br'|}\,d\sigma
  =-\oint_S \frac{\bvec{M}\times\uvec{n}}{|\br-\br'|}\,d\sigma.
\end{equation}
Inserting \eqref{eq:lec17-integrand-rewrite} and
\eqref{eq:lec17-curl-to-surface} into \eqref{eq:lec17-Am},
\begin{equation}\label{eq:lec17-Am-split}
  \Avec_M(\br)
  =\int_V \frac{\nabla'\times\bvec{M}}{|\br-\br'|}\,d^3r'
  +\oint_S \frac{\bvec{M}\times\uvec{n}}{|\br-\br'|}\,d\sigma.
\end{equation}

\subsubsection*{Reading Off the Bound Currents}
Equation \eqref{eq:lec17-Am-split} has exactly the structure of the magnetostatic
potential generated by a volume current plus a surface current,
\begin{equation}\label{eq:lec17-Am-target}
  \Avec(\br)=\frac{1}{c}\int_V \frac{\Jvec}{|\br-\br'|}\,d^3r'
  +\frac{1}{c}\oint_S \frac{\bvec{K}}{|\br-\br'|}\,d\sigma.
\end{equation}
Because the equality must hold for an arbitrary body of arbitrary shape and arbitrary
$\bvec{M}$ (and because, with no boundary, the surface piece drops and the two volume
integrals must match term by term), we identify
\begin{equation}\label{eq:lec17-bound-currents}
  \boxed{\Jvec_M=c\,\nabla\times\bvec{M},
  \qquad
  \bvec{K}_M=c\,\bvec{M}\times\uvec{n}.}
\end{equation}
\textit{Significance:} Magnetization is \emph{exactly} equivalent to a real-looking
bound current -- a volume current wherever $\bvec{M}$ curls, and a sheet current
$c\,\bvec{M}\times\uvec{n}$ wrapping the surface where $\bvec{M}$ stops abruptly.  This
is the magnetic counterpart of bound charge $\rho_b=-\nabla\cdot\bvec{P}$ in dielectrics.

\nt{The surface current is the limit of a thin volume layer: writing a surface current as
a volume current times a $\delta$-function collapses one integration and leaves the
surface integral, which is why both pieces appear together.}

%------------------------------------------------------------
\subsection{Amp\`ere's Law in Matter and the Magnetic Intensity $\Hvec$}\label{subsec:lec17-H-field}

\subsubsection*{Conceptual Background}
With the bound current known, we return to Maxwell's equations in the presence of
matter.  In magnetostatics these are $\nabla\cdot\bB=0$ and Amp\`ere's law with the
\emph{total} current.  Splitting the total current and substituting
$\Jvec_M=c\,\nabla\times\bvec{M}$ lets us absorb the matter response into a redefined
field.

\subsubsection*{Derivation}
Amp\`ere's law with the total current reads
\begin{equation}\label{eq:lec17-ampere-total}
  \nabla\times\bB=\frac{4\pi}{c}\,\Jvec_{\rm total}
  =\frac{4\pi}{c}\,\Jvec_{\rm f}+\frac{4\pi}{c}\,\Jvec_M.
\end{equation}
Insert $\Jvec_M=c\,\nabla\times\bvec{M}$; the factors of $c$ cancel in the last term,
\begin{equation}\label{eq:lec17-ampere-sub}
  \nabla\times\bB=\frac{4\pi}{c}\,\Jvec_{\rm f}+4\pi\,\nabla\times\bvec{M}.
\end{equation}
Move the magnetization term to the left,
\begin{equation}\label{eq:lec17-ampere-rearranged}
  \nabla\times\bigl(\bB-4\pi\bvec{M}\bigr)=\frac{4\pi}{c}\,\Jvec_{\rm f}.
\end{equation}
The combination $\bB-4\pi\bvec{M}$ is sourced only by free currents, so it deserves its
own name.

\dfn[dfn:lec17-H]{Magnetic Intensity $\Hvec$}{
  \begin{equation}\label{eq:lec17-H-def}
    \boxed{\Hvec=\bB-4\pi\bvec{M}.}
  \end{equation}
  Outside matter $\bvec{M}=0$, so $\Hvec=\bB$ there; only inside a magnetized body do
  $\Hvec$ and $\bB$ differ.
}
\textit{Significance:} $\Hvec$ is the magnetic field ``with the matter response divided
out''; it plays the role $\bB$ plays in vacuum but is tied to the controllable free
current.

In terms of $\Hvec$, Amp\`ere's law inside matter becomes
\begin{equation}\label{eq:lec17-curl-H}
  \boxed{\nabla\times\Hvec=\frac{4\pi}{c}\,\Jvec_{\rm f}.}
\end{equation}
\textit{Significance:} The curl of $\Hvec$ knows nothing about the magnetization -- it
responds to free currents alone, which is precisely why $\Hvec$ is the convenient field
for solving boundary problems with given conduction currents.

Taking the divergence of \eqref{eq:lec17-H-def} and using $\nabla\cdot\bB=0$
(which is exact, since there are no magnetic monopoles),
\begin{equation}\label{eq:lec17-div-H}
  \boxed{\nabla\cdot\Hvec=-4\pi\,\nabla\cdot\bvec{M}.}
\end{equation}
\textit{Significance:} Unlike $\bB$, the field $\Hvec$ is generally \emph{not}
divergence-free: wherever the magnetization has nonzero divergence, $\Hvec$ has an
effective source.  This is the seed of the magnetic scalar potential.

\begin{remark}
  In general the magnetization is itself a function of the field, $\bvec{M}=\bvec{M}(\bB)$,
  determined self-consistently by the material.  For linear media (paramagnets,
  diamagnets) the relation is linear; for ferromagnets it is complicated and hysteretic.
  Throughout this lecture we assume $\bvec{M}$ (equivalently $\Jvec_M$) is \emph{given} --
  e.g.\ a permanent magnet whose moments are fixed and only weakly disturbed by the
  applied field.
\end{remark}

%------------------------------------------------------------
\subsection{Magnetic Scalar Potential and the Electrostatic Analogy}\label{subsec:lec17-scalar-potential}

\subsubsection*{Physical Motivation}
Consider a permanent magnet sitting in a region with \emph{no free currents},
$\Jvec_{\rm f}=0$.  Then \eqref{eq:lec17-curl-H} gives $\nabla\times\Hvec=0$, exactly
like the electrostatic condition $\nabla\times\bE=0$.  A curl-free field is the gradient
of a potential, so we expect a \emph{magnetic scalar potential} and an entire
electrostatics-style solution method.

\subsubsection*{The Poisson Equation for $\psi_M$}
Since $\nabla\times\Hvec=0$, write
\begin{equation}\label{eq:lec17-H-grad}
  \Hvec=-\nabla\psi_M.
\end{equation}
Substituting into \eqref{eq:lec17-div-H} gives a Poisson equation,
\begin{equation}\label{eq:lec17-poisson-psi}
  \boxed{\nabla^2\psi_M=4\pi\,\nabla\cdot\bvec{M}=-4\pi\,\rho_M,
  \qquad
  \rho_M\equiv-\nabla\cdot\bvec{M}.}
\end{equation}
\textit{Significance:} The magnetic problem becomes formally identical to electrostatics,
with an \emph{effective magnetic charge density} $\rho_M=-\nabla\cdot\bvec{M}$.  There are
no real magnetic charges; $\rho_M$ is a mathematical bookkeeping device that lets us
import every electrostatic technique (images, separation of variables, multipoles).

\subsubsection*{Explicit Potential from the Single-Dipole Field}
To find $\psi_M$ concretely, rewrite the dipole field of Lecture 16 as a gradient.  For
a moment $\bmu$ at the origin, the full field (including the contact term at the source)
is
\begin{equation}\label{eq:lec17-Bdip-full}
  \bB_{\rm dip}(\br)
  =\frac{3\,\hat{\br}(\bmu\cdot\hat{\br})-\bmu}{r^3}
  +4\pi\,\bmu\,\delta^{(3)}(\br).
\end{equation}
We claim this equals $-\nabla\bigl(\bmu\cdot\br/r^3\bigr)+4\pi\bmu\,\delta^{(3)}(\br)$.
Indeed, differentiating the scalar $\bmu\cdot\br/r^3=\mu_j r_j r^{-3}$,
\begin{align}
  \partial_i\!\left(\frac{\mu_j r_j}{r^3}\right)
  &=\frac{\mu_i}{r^3}+\mu_j r_j\,\partial_i r^{-3} \notag\\
  &=\frac{\mu_i}{r^3}-\frac{3\,r_i(\bmu\cdot\br)}{r^5},
  \label{eq:lec17-grad-mur}
\end{align}
so that, away from the origin,
\begin{equation}\label{eq:lec17-grad-equals-field}
  -\nabla\!\left(\frac{\bmu\cdot\br}{r^3}\right)
  =\frac{3\,\hat{\br}(\bmu\cdot\hat{\br})-\bmu}{r^3},
\end{equation}
which is exactly the $r\neq0$ part of \eqref{eq:lec17-Bdip-full}.  Hence
\begin{equation}\label{eq:lec17-Bdip-gradient}
  \boxed{\bB_{\rm dip}(\br)
  =-\nabla\!\left(\frac{\bmu\cdot\br}{r^3}\right)+4\pi\,\bmu\,\delta^{(3)}(\br).}
\end{equation}
\textit{Significance:} The dipole field is (almost everywhere) a pure gradient; the
$\delta$-function carries the difference between the curl-based and gradient-based
representations and is exactly the term that reconstructs $4\pi\bvec{M}$ in the bulk.

\subsubsection*{Superposition over the Body}
Replace $\bmu\to\bvec{M}(\br')\,d^3r'$ and place each dipole at $\br'$ (so
$\br\to\br-\br'$), then integrate.  The first term integrates into a gradient and the
$\delta$-function collapses the second integral:
\begin{equation}\label{eq:lec17-BM}
  \bB_M(\br)
  =-\nabla\,\psi_M(\br)+4\pi\,\bvec{M}(\br),
  \qquad
  \psi_M(\br)=\int d^3r'\,
  \frac{\bvec{M}(\br')\cdot(\br-\br')}{|\br-\br'|^3}.
\end{equation}
Forming $\Hvec_M=\bB_M-4\pi\bvec{M}$ cancels the contact term and confirms
\begin{equation}\label{eq:lec17-HM-grad}
  \boxed{\Hvec_M=-\nabla\,\psi_M,
  \qquad
  \psi_M(\br)=\int d^3r'\,\frac{\bvec{M}(\br')\cdot(\br-\br')}{|\br-\br'|^3}.}
\end{equation}
\textit{Significance:} This is fully consistent with the abstract result
$\Hvec=-\nabla\psi_M$: starting from the explicit dipole superposition we recover both
the scalar potential and the relation $\bB=\Hvec+4\pi\bvec{M}$.  The magnetic problem
with no free currents is solved exactly like an electrostatics problem.

\subsubsection*{The Electrostatic Analogy}
The magnetic scalar potential is the literal image of the electric dipole potential.
Compare the two theories side by side:

\begin{center}
\renewcommand{\arraystretch}{1.4}
\begin{tabular}{lll}
  \toprule
  & \textbf{Electrostatics} & \textbf{Magnetostatics ($\Jvec_{\rm f}=0$)} \\
  \midrule
  Field from potential & $\bE=-\nabla\phi$ & $\Hvec=-\nabla\psi_M$ \\
  Dipole potential     & $\phi=\dfrac{\bvec{p}\cdot\br}{r^3}$
                       & $\psi_M=\dfrac{\bmu\cdot\br}{r^3}$ \\
  Source density       & $\rho$ & $\rho_M=-\nabla\cdot\bvec{M}$ \\
  Poisson equation     & $\nabla^2\phi=-4\pi\rho$ & $\nabla^2\psi_M=-4\pi\rho_M$ \\
  \bottomrule
\end{tabular}
\end{center}

\textit{Significance:} Polarization-like reasoning carries over completely:
$\bvec{M}$ plays the role of polarization, $-\nabla\cdot\bvec{M}$ the role of bound
charge, and $\psi_M$ the role of the electrostatic potential.  The next lecture exploits
this to solve concrete magnet geometries with electrostatic methods.

%------------------------------------------------------------
\subsection{Summary}\label{subsec:lec17-summary}

\begin{itemize}
  \item The total current splits as $\Jvec_{\rm total}=\Jvec_{\rm f}+\Jvec_M$, with the
        magnetization current arising microscopically from orbital motion plus an
        effective spin current; classically only $\bvec{M}=\bvec{M}_o+\bvec{M}_s$ matters.
  \item Superposing dipole vector potentials gives
        $\Avec_M=\int \bvec{M}\times(\br-\br')/|\br-\br'|^3\,d^3r'$.
  \item A vector identity converts this into bound currents
        $\Jvec_M=c\,\nabla\times\bvec{M}$ (volume) and
        $\bvec{K}_M=c\,\bvec{M}\times\uvec{n}$ (surface).
  \item Amp\`ere's law then forces $\Hvec=\bB-4\pi\bvec{M}$ with
        $\nabla\times\Hvec=(4\pi/c)\Jvec_{\rm f}$ and
        $\nabla\cdot\Hvec=-4\pi\nabla\cdot\bvec{M}$.
  \item With no free currents, $\Hvec=-\nabla\psi_M$ and
        $\nabla^2\psi_M=-4\pi\rho_M$ with effective magnetic charge
        $\rho_M=-\nabla\cdot\bvec{M}$.
  \item Explicitly $\psi_M=\int \bvec{M}(\br')\cdot(\br-\br')/|\br-\br'|^3\,d^3r'$, the
        exact analogue of the electric dipole potential -- magnetostatics in matter
        reduces to an electrostatics problem.
\end{itemize}
